% GENERATED by split_experiments.py from experiments.tex and supplementary.tex --
% do not edit directly. Appendix F is the two Experiments subsections that do not fit
% the four-page main text; Appendix G is the per-model tables.
%
\subsection{F Additional Experiments}
\subsubsection{F1 Are the baselines handicapped?}
\label{sec:fairness}

Two objections to Table~\ref{tab:main} deserve measurements rather than argument: that the
baselines are evaluated away from their published setting, and that the matched budget starves
them.

\paragraph{Released weights.}
Only Slot Attention ships trained weights; the AKOrN, CAE and SynCx repositories release data but
no checkpoints. Evaluated directly, at roughly $320\times$ our exposure budget, its released model
scores $.3747/.2219/.1165$ on CLEVR, $.4742/.2381/.2355$ on Tetrominoes and $.3205/.0980/.1204$
on Multi-dSprites. We remain ahead in all nine cells, by $0.377$ to $0.720$, so the comparison
does not rest on our retraining of it. The extra data also helps it far more \emph{off} its
training distribution than on it --- worth $+.0875$ \textsc{fg-ari} on CLEVR but $+.3003$ on
Tetrominoes and $+.1275$ on Multi-dSprites --- so a data-efficiency conclusion drawn from an
in-domain column alone would not transfer.

\paragraph{Budget.}
Comparing each baseline against its own published number, by its own evaluation code, isolates
the effect of the shortened schedule. Two of the six are \emph{not} budget-limited: AKOrN (conv)
reaches $60.5$ at $1/210$ of its schedule against a published $59.0\pm4.3$, and CAE reaches $10.3$
at $1/64$ against a published $4.0$. For those two the margin in Table~\ref{tab:main} is not an
artefact of a short schedule. The other four are limited --- AKOrN (attn) $60.4$ against
$91.0\pm0.5$, ItrSA $58.1$ against $82.5\pm1.7$, CAE\texttt{++} $16.7$ against $27.0$, and CtCAE
$12.0$ against $54.0$ --- and against those the result should be read as a data-efficiency claim.
We state both readings rather than only the favourable one. Full figures are in Appendix~G.

\subsubsection{F2 Rollout length and data scaling}
\label{sec:timestep}

The readout integrates spiking activity over $T$ simulation steps and reads the second half of the
window. Table~\ref{tab:timestep} varies $T$ with the model fixed.

\begin{table}[!tb]
\centering
\small
\begin{tabular}{r ccc r}
\toprule
$T$ & \textsc{fg-ari} & \textsc{fg-iou} & \textsc{obj-iou} & $\Delta$\,\textsc{fg-ari} \\
\midrule
$1024$ & .8421 & .7545 & .6410 & --- \\
$512$  & .8402 & \textbf{.7594} & .6340 & $-.0019$ \\
$256$  & .8159 & .7358 & .6084 & $-.0262$ \\
$128$  & .5796 & .4988 & .3985 & $-.2625$ \\
\bottomrule
\end{tabular}
\caption{CLEVR, model fixed, readout window $T/2$. Halving the rollout is free; a quarter costs
$0.02$--$0.03$; an eighth fails. This table and the scaling figures below are reported
\emph{without} the minimum-region-size rule, because both are within-model comparisons at a
fixed readout and the rule would shift every row by the same small amount.}
\label{tab:timestep}
\end{table}

Halving the rollout costs $0.0019$ \textsc{fg-ari} --- below any difference we can resolve at one
seed --- and \emph{improves} foreground IoU. Performance degrades gently to $T=256$ and collapses
at $T=128$, where the model also under-segments, predicting a foreground fraction of $0.1148$
against a truth of $0.2115$. We therefore use $T=512$ as the default operating point and $T=1024$
for the reported numbers. Shorter windows also fragment more: mean predicted regions per image
rise from $8.0$ at $T=1024$ to $13.3$ at $T=256$, against $6.36$ true objects.

\paragraph{Scaling with training data.}
Continuing training on additional unseen CLEVR images, in steps of $1{,}000$ from $10{,}000$ to
$18{,}000$, separates the three metrics. \textsc{fg-ari} runs
$.8421$, $.8426$, $.8449$, $.8460$, $.8469$, $.8459$, $.8456$, $.8470$, $.8461$: it gains
$.0028$ over the first two steps and then stays inside a band of $.0021$ for the remaining six,
with no trend. Foreground IoU runs $.7545$ to $.7755$ and matched-object IoU $.6410$ to $.6579$,
both still rising at $18{,}000$ and neither turning. More in-domain data therefore stops improving
\emph{how} patches are grouped at around $12{,}000$ images while it is still improving \emph{where
the boundaries fall}. Mean predicted regions hold at $7.7$--$8.0$ throughout, against $6.36$ true
objects, so the extra data sharpens the partition without changing how much it splits. The
\textsc{fg-ari} band is one seed wide and we read it as flat rather than converged; the two IoU
curves have not been run far enough to find their own ceiling.

\paragraph{In-domain data is not the most useful data.}
The same $5{,}000$ additional images buy more than twice as much when they come from a different
dataset. Continuing on CLEVR raises CLEVR \textsc{fg-ari} by $+.0038$; continuing on Tetrominoes
instead raises it by $+.0084$, a $2.2\times$ larger gain at an identical budget. Together with
the transfer result in Table~\ref{tab:main}, this argues that the model is not fitting a domain
but learning a grouping rule that the extra renders refine.

\subsubsection{F3 Spike timing, measured}
\label{sec:timing}

The two measurements summarised in Section 3.2, in full. Both run on the first $40$ images of
the CLEVR reserved split, on the same rollouts, so they are directly comparable.

\paragraph{Validity of the ablation.}
The affinity $Q$ is recomputed here from \texttt{component\_spikes} rather than taken from the
model, so that a modified spike trace can be fed through the same arithmetic. On the unshifted
control the recomputed $Q$ agrees with the model's own to $6.7\times10^{-6}$ maximum absolute
difference over all $4096^2$ entries, so the ablations differ from the control in exactly one
thing.

\paragraph{Separability.}
Ground-truth labels are used only to sort patch pairs into three groups and never to set a
threshold or a count. Averaged over the $40$ images, the affinity is $0.891\pm0.111$ for pairs
inside one object, $0.047\pm0.048$ for pairs in different objects, and $0.099\pm0.106$ for
object--background pairs. The within-object value is nineteen times the across-object value, and
the across-object value sits near the floor of the measure.

\paragraph{Per-patch circular shift.}
Each patch's four component traces are rolled by one offset drawn uniformly over the readout
window, independently per patch. Rate, waveform and within-patch structure survive; only the
alignment between patches is destroyed. The graph, the foreground threshold, the component
floor, the background heuristic and the greedy merge are untouched.

\begin{table}[!tb]
\centering
\small
\begin{tabular}{l ccc r}
\toprule
& \textsc{fg-ari} & \textsc{fg-iou} & \textsc{obj-iou} & regions \\
\midrule
control            & .8648 & .8028 & .6925 & 7.8 \\
per-patch shift    & .0735 & .7358 & .0325 & 562.3 \\
gates, not spikes  & .8648 & .8028 & .6925 & 7.8 \\
\bottomrule
\end{tabular}
\caption{First $40$ CLEVR reserved images. Destroying cross-patch timing collapses the
partition (\textsc{fg-ari}, \textsc{obj-iou}) while barely moving foreground detection, which
is rate-based. Substituting the phase gates for the spikes changes nothing at all.}
\label{tab:timing}
\end{table}

\paragraph{Gates in place of spikes.}
The third row replaces \texttt{component\_spikes} with \texttt{component\_gates} and changes
nothing else. Every mask is reproduced exactly, on all $40$ images and to every digit of all
three metrics. The reason is visible in the traces: with $s_{ia}=g_ie_{ia}$, the firing
indicator $e$ equals $1$ in $99.992\%$ of settled steps, the correlation between $s$ and $g$ is
$0.9997$, and the $0.008\%$ of steps where the neuron is silent never change a thresholded
decision. The timing that the shift ablation shows to be necessary is therefore the timing of
the oscillatory gate. The spiking layer carries it without loss, but the readout cannot tell
the two apart, and no claim should be made for the threshold nonlinearity on this evidence.


% Supplementary document -- AAAI-27.
%
% Holds the three tables cut from the Experiments section for the 7-page limit. Every
% number here also appears in the main paper's prose; these tables exist so a reader can
% check them laid out rather than read them in a sentence.
%
% Same preamble rules as the main paper. This file is \input by main_supplementary.tex,
% which compiles to a SEPARATE PDF -- AAAI-27 takes the supplementary as its own upload,
% not as extra pages of the paper. It is deliberately not \input by main_aaai.tex.

\subsection{G Supplementary Tables}

\subsubsection{G1 Released weights}

These figures are given in prose in the main paper.

% table*, not table: three columns of metric triples overflow a single 3.3in column by
% 122pt. table* spans both columns of the supplementary, which is also two-column.
\begin{table*}[!tb]
\centering
\small
\setlength{\tabcolsep}{4pt}
\begin{tabular}{l ccc}
\toprule
& CLEVR & Tetrominoes & Multi-dSprites \\
\midrule
Slot Attention, matched budget   & .2872\,/\,.2220\,/\,.0801 & .1739\,/\,.2192\,/\,.1071 & .1930\,/\,.0876\,/\,.0680 \\
Slot Attention, released weights & .3747\,/\,.2219\,/\,.1165 & .4742\,/\,.2381\,/\,.2355 & .3205\,/\,.0980\,/\,.1204 \\
\textbf{Ours, matched budget}    & \textbf{.8421\,/\,.7545\,/\,.6410} & \textbf{.8580\,/\,.7885\,/\,.8173} & \textbf{.6974\,/\,.8177\,/\,.7628} \\
\bottomrule
\end{tabular}
\caption{Released weights ($\approx\!32$M exposures) against the matched-budget retraining
($100{,}032$ exposures) and against our model. Triples are
\textsc{ari}\,/\,\textsc{f-iou}\,/\,\textsc{o-iou}. All three rows are reported
\emph{without} the minimum-region-size rule, so the released and retrained Slot Attention
rows are compared on identical terms; Table~1 of the main paper applies the rule to every
model and our row there reads $.8408/.7528/.6387$, $.8568/.7950/.8170$ and
$.6966/.8190/.7625$.}
\label{tab:s-released}
\end{table*}

\subsubsection{G2 Training budget per baseline}

\begin{table}[!tb]
\centering
\small
\begin{tabular}{l c c c}
\toprule
Model & Budget & Ours & Published \\
\midrule
AKOrN (conv) & $1/210$ & \textbf{60.5} & $59.0\pm4.3$ \\
CAE          & $1/64$  & \textbf{10.3} & $4.0$ \\
\midrule
AKOrN (attn) & $1/210$ & 60.4 & $91.0\pm0.5$ \\
ItrSA        & $1/210$ & 58.1 & $82.5\pm1.7$ \\
CAE\texttt{++} & $1/64$ & 16.7 & $27.0$ \\
CtCAE        & $1/64$  & 12.0 & $54.0$ \\
\bottomrule
\end{tabular}
\caption{Each baseline at our budget against its own published figure, by its own metric and
evaluation code. \emph{Budget} is the fraction of the published training schedule our
matched-exposure protocol gives it. The top block is \emph{not} budget-limited: AKOrN (conv)
exceeds its published CLEVR number at $1/210$ of its schedule and CAE exceeds its own at $1/64$.
The bottom block is.}
\label{tab:s-budget}
\end{table}

\subsubsection{G3 Scaling with training data}

\begin{table}[!tb]
\centering
\small
\begin{tabular}{r ccc r}
\toprule
Images & \textsc{fg-ari} & \textsc{fg-iou} & \textsc{obj-iou} & $\Delta$\,\textsc{fg-ari} \\
\midrule
$10{,}000$ & .8421 & .7545 & .6410 & --- \\
$11{,}000$ & .8426 & .7539 & .6427 & $+.0005$ \\
$12{,}000$ & \textbf{.8449} & \textbf{.7607} & \textbf{.6461} & $+.0028$ \\
\bottomrule
\end{tabular}
\caption{\emph{Images} is the cumulative count of unique training images. Each row adds $1{,}000$
previously unseen CLEVR images for $10$ epochs, resuming from the previous row. CLEVR
evaluation.}
\label{tab:s-scaling}
\end{table}
